% ============================================================
%  DOCUMENTO 1: REVIEW (estilo articulo, formato SIMBIG - sin indice)
%  Reconocimiento de voz, traduccion automatica y sintesis de voz aplicados
%  al quechua para la comunicacion medico-paciente: revision sistematica
%  Compilar:  pdflatex review -> biber review -> pdflatex x2
% ============================================================
\documentclass[12pt,a4paper]{article}
% ============================================================
%  Preambulo compartido - APA7 + espanol + graficos
%  Se incluye con % ============================================================
%  Preambulo compartido - APA7 + espanol + graficos
%  Se incluye con % ============================================================
%  Preambulo compartido - APA7 + espanol + graficos
%  Se incluye con \input{../preambulo.tex} en cada documento.
%  Compilar con: pdflatex -> biber -> pdflatex -> pdflatex
% ============================================================

% ---- Codificacion e idioma ----
\usepackage[utf8]{inputenc}
\usepackage[T1]{fontenc}
\usepackage{lmodern}
\usepackage[spanish,es-tabla]{babel}
\usepackage{csquotes}

% ---- Geometria (margenes tipo tesis; ajustar a SINVIC) ----
\usepackage[a4paper,left=3cm,right=2.5cm,top=2.5cm,bottom=2.5cm]{geometry}

% ---- Interlineado 1.5 ----
\usepackage{setspace}
\onehalfspacing

% ---- Tipografia de titulos ----
\usepackage{titlesec}
\titleformat{\section}{\normalfont\Large\bfseries}{\thesection}{1em}{}
\titleformat{\subsection}{\normalfont\large\bfseries}{\thesubsection}{1em}{}

% ---- Tablas y figuras ----
\usepackage{graphicx}
\usepackage{booktabs}
\usepackage{array}
\usepackage{multirow}
\usepackage{tabularx}
\usepackage{float}
\usepackage[font=small,labelfont=bf]{caption}

% ---- Graficos vectoriales (pgfplots / tikz) ----
\usepackage{pgfplots}
\pgfplotsset{compat=1.18}
\usepackage{tikz}
\usetikzlibrary{shapes.geometric,arrows.meta,positioning,calc}

% ---- Colores ----
\usepackage{xcolor}
\definecolor{azulUNSAAC}{RGB}{0,71,133}
\definecolor{grisSuave}{RGB}{240,240,240}

% ---- Bibliografia APA7 (biblatex + biber) ----
\usepackage[
  backend=biber,
  style=apa,
  sortcites=true,
  sorting=nyt,
  language=spanish
]{biblatex}
\DeclareLanguageMapping{spanish}{spanish-apa}
\addbibresource{../referencias.bib}

% ---- Listas e indices (algoritmos, ecuaciones, anexos; tablas largas) ----
% IMPORTANTE: tocloft debe cargarse ANTES de hyperref.
\usepackage{longtable}
\usepackage{algorithm}
\usepackage{algpseudocode}
\usepackage{tocloft}

% ---- Hipervinculos ----
\usepackage[hidelinks]{hyperref}
\hypersetup{
  colorlinks=true,
  linkcolor=black,
  citecolor=azulUNSAAC,
  urlcolor=azulUNSAAC
}

% ---- Icono ORCID (dibujado con tikz, compatible con pdflatex) ----
\usepackage{orcidlink}

% ---- Encabezado / pie de pagina ----
\usepackage{fancyhdr}
\pagestyle{fancy}
\fancyhf{}
\fancyfoot[C]{\thepage}
\renewcommand{\headrulewidth}{0pt}

% ---- Sangria de primer parrafo ----
\usepackage{indentfirst}
 en cada documento.
%  Compilar con: pdflatex -> biber -> pdflatex -> pdflatex
% ============================================================

% ---- Codificacion e idioma ----
\usepackage[utf8]{inputenc}
\usepackage[T1]{fontenc}
\usepackage{lmodern}
\usepackage[spanish,es-tabla]{babel}
\usepackage{csquotes}

% ---- Geometria (margenes tipo tesis; ajustar a SINVIC) ----
\usepackage[a4paper,left=3cm,right=2.5cm,top=2.5cm,bottom=2.5cm]{geometry}

% ---- Interlineado 1.5 ----
\usepackage{setspace}
\onehalfspacing

% ---- Tipografia de titulos ----
\usepackage{titlesec}
\titleformat{\section}{\normalfont\Large\bfseries}{\thesection}{1em}{}
\titleformat{\subsection}{\normalfont\large\bfseries}{\thesubsection}{1em}{}

% ---- Tablas y figuras ----
\usepackage{graphicx}
\usepackage{booktabs}
\usepackage{array}
\usepackage{multirow}
\usepackage{tabularx}
\usepackage{float}
\usepackage[font=small,labelfont=bf]{caption}

% ---- Graficos vectoriales (pgfplots / tikz) ----
\usepackage{pgfplots}
\pgfplotsset{compat=1.18}
\usepackage{tikz}
\usetikzlibrary{shapes.geometric,arrows.meta,positioning,calc}

% ---- Colores ----
\usepackage{xcolor}
\definecolor{azulUNSAAC}{RGB}{0,71,133}
\definecolor{grisSuave}{RGB}{240,240,240}

% ---- Bibliografia APA7 (biblatex + biber) ----
\usepackage[
  backend=biber,
  style=apa,
  sortcites=true,
  sorting=nyt,
  language=spanish
]{biblatex}
\DeclareLanguageMapping{spanish}{spanish-apa}
\addbibresource{../referencias.bib}

% ---- Listas e indices (algoritmos, ecuaciones, anexos; tablas largas) ----
% IMPORTANTE: tocloft debe cargarse ANTES de hyperref.
\usepackage{longtable}
\usepackage{algorithm}
\usepackage{algpseudocode}
\usepackage{tocloft}

% ---- Hipervinculos ----
\usepackage[hidelinks]{hyperref}
\hypersetup{
  colorlinks=true,
  linkcolor=black,
  citecolor=azulUNSAAC,
  urlcolor=azulUNSAAC
}

% ---- Icono ORCID (dibujado con tikz, compatible con pdflatex) ----
\usepackage{orcidlink}

% ---- Encabezado / pie de pagina ----
\usepackage{fancyhdr}
\pagestyle{fancy}
\fancyhf{}
\fancyfoot[C]{\thepage}
\renewcommand{\headrulewidth}{0pt}

% ---- Sangria de primer parrafo ----
\usepackage{indentfirst}
 en cada documento.
%  Compilar con: pdflatex -> biber -> pdflatex -> pdflatex
% ============================================================

% ---- Codificacion e idioma ----
\usepackage[utf8]{inputenc}
\usepackage[T1]{fontenc}
\usepackage{lmodern}
\usepackage[spanish,es-tabla]{babel}
\usepackage{csquotes}

% ---- Geometria (margenes tipo tesis; ajustar a SINVIC) ----
\usepackage[a4paper,left=3cm,right=2.5cm,top=2.5cm,bottom=2.5cm]{geometry}

% ---- Interlineado 1.5 ----
\usepackage{setspace}
\onehalfspacing

% ---- Tipografia de titulos ----
\usepackage{titlesec}
\titleformat{\section}{\normalfont\Large\bfseries}{\thesection}{1em}{}
\titleformat{\subsection}{\normalfont\large\bfseries}{\thesubsection}{1em}{}

% ---- Tablas y figuras ----
\usepackage{graphicx}
\usepackage{booktabs}
\usepackage{array}
\usepackage{multirow}
\usepackage{tabularx}
\usepackage{float}
\usepackage[font=small,labelfont=bf]{caption}

% ---- Graficos vectoriales (pgfplots / tikz) ----
\usepackage{pgfplots}
\pgfplotsset{compat=1.18}
\usepackage{tikz}
\usetikzlibrary{shapes.geometric,arrows.meta,positioning,calc}

% ---- Colores ----
\usepackage{xcolor}
\definecolor{azulUNSAAC}{RGB}{0,71,133}
\definecolor{grisSuave}{RGB}{240,240,240}

% ---- Bibliografia APA7 (biblatex + biber) ----
\usepackage[
  backend=biber,
  style=apa,
  sortcites=true,
  sorting=nyt,
  language=spanish
]{biblatex}
\DeclareLanguageMapping{spanish}{spanish-apa}
\addbibresource{../referencias.bib}

% ---- Listas e indices (algoritmos, ecuaciones, anexos; tablas largas) ----
% IMPORTANTE: tocloft debe cargarse ANTES de hyperref.
\usepackage{longtable}
\usepackage{algorithm}
\usepackage{algpseudocode}
\usepackage{tocloft}

% ---- Hipervinculos ----
\usepackage[hidelinks]{hyperref}
\hypersetup{
  colorlinks=true,
  linkcolor=black,
  citecolor=azulUNSAAC,
  urlcolor=azulUNSAAC
}

% ---- Icono ORCID (dibujado con tikz, compatible con pdflatex) ----
\usepackage{orcidlink}

% ---- Encabezado / pie de pagina ----
\usepackage{fancyhdr}
\pagestyle{fancy}
\fancyhf{}
\fancyfoot[C]{\thepage}
\renewcommand{\headrulewidth}{0pt}

% ---- Sangria de primer parrafo ----
\usepackage{indentfirst}


\title{\vspace{-1.2cm}\bfseries Reconocimiento de voz, traducci{\'o}n autom{\'a}tica y s{\'i}ntesis de voz
aplicados al quechua para la comunicaci{\'o}n m{\'e}dico-paciente: una revisi{\'o}n sistem{\'a}tica de literatura%
\thanks{Trabajo del curso Seminario de Tesis II. Docente: Mtro.\ Ing.\ Ra{\'u}l Huillca Huallparimachi.}}
\author{Alain Armando Condori Contreras\,(021501)\\[2pt]
  \normalsize Escuela Profesional de Ingenier{\'i}a Inform{\'a}tica y de Sistemas\\[1pt]
  \normalsize Universidad Nacional de San Antonio Abad del Cusco\\[1pt]
  \normalsize \texttt{021501@unsaac.edu.pe}}
\date{Cusco, Per{\'u} --- 2026}

\begin{document}
\maketitle

% ------------------- RESUMEN -------------------
\begin{abstract}
\noindent Esta revisi{\'o}n sintetiza el estado del arte sobre traducci{\'o}n autom{\'a}tica del par
Quechua--Espa{\~n}ol y las tecnolog{\'i}as de reconocimiento de voz (ASR) y s{\'i}ntesis de voz
(TTS) aplicables a lenguas de bajo recurso, con foco en el dominio de la consulta m{\'e}dica.
Se organiza en tres ejes: (i) la barrera ling{\"u}{\'i}stica en salud en contextos
quechuahablantes, (ii) los modelos preentrenados de c{\'o}digo abierto disponibles
(Whisper, MMS, NLLB) y (iii) los vac{\'i}os metodol{\'o}gicos que justifican el presente
trabajo.

\vspace{0.3cm}
\noindent\textbf{Palabras clave:} traducci{\'o}n autom{\'a}tica, quechua, PLN, lenguas de bajo recurso, ASR, TTS.
\end{abstract}

% ------------------- 1. INTRODUCCION -------------------
\section{Introducci{\'o}n}
La brecha comunicativa entre el personal de salud y los pacientes quechuahablantes es un
problema documentado y vigente en la regi{\'o}n Cusco \parencite{begazo2026}. Esta revisi{\'o}n
busca mapear qu{\'e} soluciones tecnol{\'o}gicas existen, qu{\'e} tan maduras est{\'a}n para el
quechua y d{\'o}nde se ubican los vac{\'i}os que este proyecto puede cubrir.

\subsection{Objetivo de la revisi{\'o}n}
Identificar y comparar los modelos preentrenados de ASR, traducci{\'o}n autom{\'a}tica y TTS
que soportan quechua, evaluando su viabilidad para un prototipo de traducci{\'o}n
bidireccional en consulta m{\'e}dica.

\subsection{Metodolog{\'i}a de b{\'u}squeda}
Se revisaron fuentes en Google Scholar, arXiv, ACL Anthology y repositorios
institucionales peruanos, con los t{\'e}rminos: \emph{quechua machine translation},
\emph{low-resource speech recognition}, \emph{multilingual TTS} y \emph{barrera
ling{\"u}{\'i}stica salud Per{\'u}}. Criterio de inclusi{\'o}n: publicaciones 2019--2026
con relevancia directa al par Quechua--Espa{\~n}ol o a lenguas de bajo recurso comparables.

% ------------------- 2. BARRERA LINGUISTICA -------------------
\section{La barrera ling{\"u}{\'i}stica en salud}
En Cusco, m{\'a}s del 55\,\% de la poblaci{\'o}n tiene el quechua como lengua materna
\parencite{essalud2023}. Los estudios cualitativos confirman que la desconexi{\'o}n
idiom{\'a}tica vulnera la relaci{\'o}n terap{\'e}utica \parencite{begazo2026,cuaila2022}, y a
ello se suma el analfabetismo en mujeres rurales mayores \parencite{inei2023,dispurse2023}.

La Figura~\ref{fig:barreras} resume las magnitudes de la problem{\'a}tica.

\begin{figure}[H]
\centering
\begin{tikzpicture}
\begin{axis}[
    ybar,
    width=12cm, height=7cm,
    bar width=28pt,
    ymin=0, ymax=100,
    ylabel={Porcentaje (\%)},
    symbolic x coords={Lengua materna,Asegurados EsSalud,Analfab. mujer rural},
    xtick=data,
    nodes near coords,
    nodes near coords style={font=\small},
    x tick label style={align=center,font=\small,text width=2.5cm},
    enlarge x limits=0.25,
    ymajorgrids=true,
    grid style=dashed,
]
\addplot[fill=azulUNSAAC] coordinates {
    (Lengua materna,55)
    (Asegurados EsSalud,15)
    (Analfab. mujer rural,22.8)
};
\end{axis}
\end{tikzpicture}
\caption{Magnitud de la barrera ling{\"u}{\'i}stica y de alfabetizaci{\'o}n en Cusco. Elaboraci{\'o}n propia con datos de \textcite{essalud2023} e \textcite{inei2023}.}
\label{fig:barreras}
\end{figure}

% ------------------- 3. MODELOS -------------------
\section{Modelos preentrenados de c{\'o}digo abierto}
La estrategia dominante en lenguas de bajo recurso ya no es entrenar desde cero, sino
\emph{adaptar} (fine-tuning) modelos multiling{\"u}es masivos preentrenados sobre grandes
vol{\'u}menes de datos, transfiriendo su conocimiento a la lengua objetivo con relativamente
pocos ejemplos etiquetados \parencite{pratap2024mms}. A continuaci{\'o}n se revisan los tres
componentes del pipeline.

\subsection{Reconocimiento de voz (ASR)}
Whisper \parencite{radford2023whisper} es un modelo entrenado con 680\,000 horas de audio
multiling{\"u}e mediante supervisi{\'o}n d{\'e}bil, robusto a ruido y acentos, con quechua entre
sus idiomas declarados. MMS \parencite{pratap2024mms} extiende el ASR a m{\'a}s de 1\,100
idiomas, incluyendo variantes del quechua, apoy{\'a}ndose en datos religiosos alineados. En
ambos casos, el desempe{\~n}o \emph{real} sobre vocabulario cl{\'i}nico y sobre las variantes
del quechua sure{\~n}o (Cusco--Collao) es una pregunta abierta que este proyecto debe medir
emp{\'i}ricamente, pues los benchmarks publicados no cubren el dominio m{\'e}dico.

\subsection{Traducci{\'o}n autom{\'a}tica}
NLLB \parencite{costajussa2022nllb} incluye pares con quechua, aunque con calidad
heterog{\'e}nea seg{\'u}n la variante. Para el par Quechua--Espa{\~n}ol existen antecedentes
espec{\'i}ficos: \textcite{ortega2020quechua} demostraron que la traducci{\'o}n neural en una
lengua polisint{\'e}tica como el quechua se beneficia del tratamiento morfol{\'o}gico
(segmentaci{\'o}n en subpalabras) frente a enfoques palabra a palabra. El \emph{shared task}
de AmericasNLP \parencite{mager2021americasnlp} consolid{\'o} conjuntos de datos y l{\'i}neas
base para diez lenguas ind{\'i}genas americanas, el quechua entre ellas, y evidenci{\'o} que la
escasez y el ruido de los corpus paralelos siguen siendo el cuello de botella principal.

\subsection{S{\'i}ntesis de voz (TTS)}
MMS-TTS \parencite{pratap2024mms} provee voces para cientos de idiomas; su naturalidad e
inteligibilidad en quechua es el principal riesgo t{\'e}cnico a validar, dado que la calidad
percibida de la voz es cr{\'i}tica cuando el destinatario es un paciente que no sabe leer. Por
ello el proyecto trata el TTS como componente a caracterizar antes de comprometerlo como
objetivo central.

\subsection{Recursos de datos para el quechua}
Un habilitador reciente es la aparici{\'o}n de corpus de voz para lenguas peruanas: el corpus
\emph{Huqariq} \parencite{zevallos2022huqariq} re{\'u}ne horas de habla en quechua, aimara,
ash{\'a}ninka y shipibo, y constituye una base valiosa para el fine-tuning del ASR, aunque no
est{\'a} especializado en vocabulario m{\'e}dico ---brecha que este proyecto aborda con la
construcci{\'o}n de un corpus cl{\'i}nico propio.

\subsection{M{\'e}tricas de evaluaci{\'o}n}
La comparaci{\'o}n de sistemas exige m{\'e}tricas estandarizadas. Para el ASR se emplea la
\emph{tasa de error de palabra} (WER), que mide inserciones, sustituciones y eliminaciones
frente a una transcripci{\'o}n de referencia; a menor WER, mejor. Para la traducci{\'o}n se usan
BLEU \parencite{papineni2002bleu} y chrF, que comparan la salida con traducciones humanas de
referencia. Para el TTS, la calidad se estima con la \emph{Mean Opinion Score} (MOS), un
promedio de valoraciones subjetivas de inteligibilidad y naturalidad. La
Tabla~\ref{tab:metricas} resume su uso en este proyecto.

\begin{table}[H]
\centering
\caption{M{\'e}tricas de evaluaci{\'o}n por componente del sistema.}
\label{tab:metricas}
\small
\begin{tabularx}{\linewidth}{@{}l l X@{}}
\toprule
\textbf{Componente} & \textbf{M{\'e}trica} & \textbf{Interpretaci{\'o}n} \\
\midrule
ASR         & WER          & Menor es mejor; error de transcripci{\'o}n \\
Traducci{\'o}n & BLEU / chrF & Mayor es mejor; similitud con referencia humana \\
TTS         & MOS          & 1--5; percepci{\'o}n de naturalidad e inteligibilidad \\
Global      & Tiempo de respuesta & Latencia extremo a extremo por turno \\
\bottomrule
\end{tabularx}
\end{table}

\begin{table}[H]
\centering
\caption{Comparaci{\'o}n de modelos preentrenados con soporte de quechua.}
\label{tab:modelos}
\small
\begin{tabularx}{\linewidth}{@{}l l X c@{}}
\toprule
\textbf{Modelo} & \textbf{Tarea} & \textbf{Cobertura quechua} & \textbf{Licencia} \\
\midrule
Whisper & ASR            & Declarada (multiling{\"u}e) & MIT \\
MMS     & ASR / TTS      & $>$1000 idiomas, incl. quechua & CC-BY-NC \\
NLLB    & Traducci{\'o}n & Pares con quechua & CC-BY-NC \\
\bottomrule
\end{tabularx}
\end{table}

% ------------------- 4. VACIOS -------------------
\section{Vac{\'i}os identificados}
El cruce entre la literatura del problema social \parencite{begazo2026,cuaila2022} y la
literatura t{\'e}cnica \parencite{radford2023whisper,pratap2024mms,mager2021americasnlp}
revela una desconexi{\'o}n: el problema est{\'a} bien documentado y la tecnolog{\'i}a base
existe, pero nadie ha unido ambos extremos. Espec{\'i}ficamente:
\begin{enumerate}
  \item \textbf{Ausencia de corpus cl{\'i}nico.} Los recursos existentes
        \parencite{zevallos2022huqariq} son de dominio general; no hay un corpus
        Quechua--Espa{\~n}ol de vocabulario de consulta (triaje, s{\'i}ntomas, indicaciones)
        validado con personal de salud.
  \item \textbf{Falta de evaluaci{\'o}n en dominio.} No se reportan valores de WER ni BLEU de
        los modelos preentrenados sobre habla y texto m{\'e}dico en quechua sure{\~n}o.
  \item \textbf{TTS sin caracterizar.} La inteligibilidad del TTS en quechua para usuarios con
        baja alfabetizaci{\'o}n ---justamente el caso de uso cr{\'i}tico--- no ha sido medida.
  \item \textbf{Diseño de interacci{\'o}n no estudiado.} No existe evidencia sobre qu{\'e}
        modalidad de salida (texto vs. texto+voz) reduce mejor los errores de interpretaci{\'o}n
        seg{\'u}n el tipo de usuario.
\end{enumerate}

% ------------------- 5. CONCLUSIONES -------------------
\section{Conclusiones}
La revisi{\'o}n confirma tres puntos. Primero, la barrera ling{\"u}{\'i}stica en salud es un
problema real, vigente y acad{\'e}micamente reconocido en Cusco
\parencite{begazo2026,essalud2023}. Segundo, existe una base tecnol{\'o}gica madura de modelos
preentrenados de c{\'o}digo abierto (Whisper, MMS, NLLB) y antecedentes de traducci{\'o}n neural
para el quechua \parencite{ortega2020quechua,mager2021americasnlp} que hacen viable la
adaptaci{\'o}n en lugar del entrenamiento desde cero. Tercero, persiste un vac{\'i}o claro en la
aplicaci{\'o}n de esa tecnolog{\'i}a al dominio cl{\'i}nico con evaluaci{\'o}n emp{\'i}rica. El
presente proyecto se posiciona en ese vac{\'i}o, priorizando ASR y traducci{\'o}n, midiendo con
m{\'e}tricas estandarizadas (WER, BLEU, MOS) y tratando el TTS como componente a validar dentro
del alcance del semestre.

% ------------------- BIBLIOGRAFIA -------------------
\newpage
\printbibliography[title={Referencias}]

\end{document}
